\honest{Make an Extraordinary Effort}{Eliezer Yudkowsky}{2008}

My second epigraph says: ``In important matters, a `strong' effort usually results in only mediocre results.'' I have seen this ``over and over again.''

A level beyond ``I want to become stronger'' is isshoukenmei: ``make a desperate effort,'' as if your life were at stake. The word began as the loyalty a samurai offered in return for his position, and its characters include ``life'' and ``land.'' It belonged to bushido, which was not only about fighting. One source says the variant issho kenmei means an all-out effort on a single point and isshou kenmei a lifelong one. \nb{Japanese sources agree in outline. The older form, isshokenmei, referred to the land a samurai would defend with his life, and the spelling with the long vowel developed from it. Dictionaries give both spellings the same everyday meaning, ``with all one's might.''}

I try not to praise the East too much, since the West hears only selected parts of it. But having a compact phrase for this is one point where Japan scores higher than America. A Japanese parent might say it to a student before exams, and it is not cheap hypocrisy there, because exams are taken very seriously.

Why don't self-styled rationalists do much better in life? From my own history, it takes a tremendous amount of rationality to stop making stupid mistakes. Robert Aumann proved that Bayesians with the same priors cannot agree to disagree, and he is a believing Orthodox Jew. Knowing the math is not enough to save him. It takes heuristics and biases, social and evolutionary psychology, and also isshoukenmei, a desperate effort to be rational, to rise above Aumann's level. \nb{Yudkowsky had used this example twice before. In a 2005 interview Aumann described his religion as ``mainly an emotional and aesthetic'' experience; see ``Crisis of Faith.''}

Should I peddle rationality in Japan instead? But Japan does not lead the United States in science, despite more studious students, and does not rule the world, though many expected it to in the 1980s. In the West, ``The squeaky wheel gets the grease''; in Japan, ``The nail that sticks up gets hammered down.'' Entrepreneurship, risk-taking and leaving the herd ``would seem to count for at least as much as desperate efforts.'' \nb{The evidence is one comparison between two countries, which differ in many ways besides these.}

So there is a second virtue. A desperate effort can lift more weight than usual, but lifting a truck takes something out of the ordinary: something you were not taught in school, that others do not expect and may not understand, outside your routine, for which you have no mental program, bypassing the System. ``This is not included in isshokenmei, or Japan would be a very different place.''

The second virtue is higher, and more dangerous. A desperate effort to lift a weight may tear a muscle. A creative idea gone wrong can blow up the truck and any number of bystanders. A businessman who makes a desperate effort to avoid bankruptcy differs from one who goes to extraordinary lengths to hide an embezzlement.

A friend of my little brother's once wanted to change a game's rules just because playing by ordinary rules was boring. I told him: ``Don't violate rules for the sake of violating them. If you break the rules only when you have an overwhelmingly good reason to do so, you will have more than enough trouble to last you the rest of your life.''

Still, we could value this virtue more. People tell me that Friendly AI is futile because corporations that care only for profit, or the military, will build the first AIs. Does it occur to them to try for something other than the default outcome? If I believed what they believe, I would not shrug and go on my way.

Others tell me to get degrees and publish ordinary papers first, ``at least a ten-year detour'' to do everything the default way. Do they think humanity can survive if every single person does everything the ordinary way? I do not make plans that need a majority, or even 10\%, of people to leave their comfort zone. So I plan for a small, privately funded project, a ``brain in a box in a basement.'' Funding it needs only a tiny fraction of six billion people to think for more than five seconds about a question that is not prepackaged. There is an awful justice in that: ``the life or death of the human species'' depends on a few people doing something a little extraordinary. The penalty is disproportionate, but most of Nature's challenges have no justice at all. Then I leave ``the details of that debate'' aside, still stunned at how often a single extraordinary element is taken as an absolute obstacle. \nb{The premise is stated here, not argued. Whether the extraordinary course was right in these two cases, both from the author's own work, depends on it.}

Keeping things ordinary can be a useful heuristic, I grant, and the risks of the extraordinary accumulate. But ordinariness has a cost too, and it is not always one you can afford. People imagine futures that will not be much fun, or that have a tinge of sadness and loss, and do not ask whether we could do better, because the sadness seems ordinary. As a smiling man once said, ``It's all part of the plan.'' \nb{Probably the Joker in \textsc{The Dark Knight} (2008), who says that nobody panics when things go ``according to plan,'' even when the plan is horrifying.}
